\chapter*{Preface}
\addcontentsline{toc}{chapter}{Preface}

This book grew out of a conviction that the two great traditions of reinforcement learning\,---\,the
classical mathematical theory and the modern practice of aligning large language
models\,---\,deserve to be told as a single, coherent story.  In recent years, ideas that spent
decades confined to the pages of control theory and dynamic programming have burst into the
centre of AI development, powering the training of systems that write, reason, and converse.
That journey is the subject of the present volume.

\medskip

\textbf{A note on inspiration and originality.}\quad
Every textbook builds on earlier work, and this one is no exception.  We wish to
acknowledge two books that have shaped the broader field and, inevitably, our
own perspective---while clarifying what distinguishes the present volume.

The first is \emph{Reinforcement Learning: An Introduction} by Richard
S.\ Sutton and Andrew G.\ Barto, the standard reference for classical RL\@.
Sutton and Barto established the pedagogical progression---from bandits
through MDPs to temporal-difference methods---that has guided virtually every
subsequent course.  Our treatment in Parts~I and~II follows the same logical
arc, but departs in emphasis and presentation: we include tighter convergence
analyses informed by recent results, introduce Generalised Advantage
Estimation and PPO within the core narrative rather than as afterthoughts, and
motivate every classical idea with an eye toward its eventual role in
language-model training.  These choices reflect lessons learned from designing
and delivering university-level RL courses as well as from applying these
algorithms in production settings at scale.

The second is \emph{The RLHF Book} by Nathan Lambert, the most thorough
practical guide to reinforcement learning from human feedback for language
models yet written.  Lambert demonstrated that the algorithmic ideas behind
ChatGPT, Claude, and their successors could be explained precisely and
accessibly.  Parts~III and~IV of our book cover overlapping territory---reward
modelling, preference optimisation, PPO for LLMs---but extend the discussion
in directions shaped by our own research and engineering experience: we
provide a unified formal treatment connecting DPO, IPO, and KTO
through the KL-constrained RL objective, and cover ORPO as a distinct
reference-model-free approach; dedicate full chapters to GRPO and
reinforcement learning with verifiable rewards (RLVR), topics that
postdate Lambert's initial treatment; and include a substantial chapter on
multi-step agentic systems that bridges RL theory with the emerging practice
of tool-using language agents.

More broadly, this book is not a synthesis of the two works above but an
attempt to tell a \emph{new} story: one continuous narrative that begins with
Bellman's principle of optimality and ends with the training pipelines behind
modern reasoning models.  That narrative is informed by hands-on experience
with large-scale RLHF training, reward-model debugging, and the practical
pitfalls that no purely theoretical account can capture.  Where we succeed in
making this bridge feel natural, the credit is shared with the community whose
open research made it possible.  Where we fall short, the responsibility is
entirely our own.


\medskip

\textbf{What this book is trying to do.}\quad
The central goal of this volume is to build a single, unbroken bridge between the
mathematical foundations laid by Bellman, Sutton, and Barto and the practical methods that
now govern how large language models are trained and aligned.  The path from Bellman's
principle of optimality to a working RLHF pipeline is shorter than it might appear, but it
requires genuine fluency in both worlds.  A practitioner who understands only the engineering
recipes misses the theoretical guarantees and the diagnostic intuition that theory provides.
A theorist who has not engaged with the realities of language-model training\,---\,the
instability of reward models, the sensitivity of PPO hyperparameters, the challenge of
specifying human preferences\,---\,is ill-equipped to contribute to the most important
application of RL today.

Parts~I and~II of this book (covering bandits, MDPs, dynamic programming, Monte Carlo
methods, temporal-difference learning, function approximation, policy gradients, and
actor-critic methods) are inspired by Sutton and Barto's rigorous treatment and can be read
entirely independently as a self-contained classical RL textbook.  Parts~III and~IV build on
that foundation to cover reward modelling, RLHF, Direct Preference Optimisation, Group
Relative Policy Optimisation, multi-step agentic systems, and open research frontiers,
drawing substantially on Lambert's practical insights.

\medskip

\textbf{Who this book is for.}\quad
We have written with three overlapping audiences in mind.

\emph{Graduate students in machine learning and AI} who want a mathematically honest account
of reinforcement learning that does not stop at the classical results but carries them all the
way to the frontier of language-model alignment.

\emph{ML engineers and practitioners} who have encountered RLHF in the context of language
models and want to understand the theoretical machinery underneath the implementation, so
that they can debug, extend, and improve it with confidence.

\emph{Researchers transitioning from classical supervised or unsupervised learning to
reinforcement learning}, including those whose primary interest is alignment, safety, or
evaluation, who need a solid grounding in the underlying concepts before engaging with the
research literature.

The book is not aimed at complete beginners to machine learning; some prior exposure to the
ideas of the field is assumed.  It is, however, self-contained with respect to reinforcement
learning itself: no prior knowledge of RL is required.

\medskip

\textbf{Prerequisites.}\quad
The reader is assumed to be comfortable with basic probability theory (random variables,
expectations, conditional probability, Bayes' rule), linear algebra (vectors, matrices,
eigenvalues), and single- and multi-variable calculus (derivatives, gradients, the chain
rule).  Familiarity with the core concepts of supervised machine learning\,---\,loss
functions, gradient descent, neural networks, overfitting\,---\,will make Parts~III and~IV
significantly more accessible, though we recall the key ideas where they are needed.
Python fluency is useful for the implementation-oriented material in the appendices but is
not necessary to follow the theoretical development.

\medskip

\textbf{How to read this book.}\quad
The four parts of the book form a natural progression, but readers with different backgrounds
may wish to enter at different points.

\emph{Part~I (Foundations)} develops the classical theory from scratch: multi-armed bandits,
Markov decision processes, dynamic programming, Monte Carlo methods, and
temporal-difference learning.  A reader with no prior exposure to RL should work through
these chapters in order.

\emph{Part~II (Function Approximation and Deep RL)} moves to settings where tabular methods
are no longer feasible.  It covers linear and nonlinear value-function approximation, policy
gradient methods, and the actor-critic family, culminating in Proximal Policy Optimisation.
Together with Part~I, this material constitutes a complete introduction to classical and
deep RL that stands on its own.

\emph{Part~III (RL for Language Models, Chapters~10--14)} applies everything developed in
Parts~I and~II to the modern challenge of training and aligning large language models.
Readers who are already fluent in classical RL may proceed directly here.  These chapters
cover reward shaping and reward models (Chapter~10), the language-model training pipeline
from pre-training through SFT (Chapter~11), Direct Preference Optimisation and its
variants (Chapter~12), Group Relative Policy Optimisation and online RL with verifiable
rewards (Chapter~13), and the practical engineering of RLHF at scale (Chapter~14).

\emph{Part~IV (Frontiers, Chapters~15--17)} explores the cutting edge of the field:
reasoning, tool use, and reinforcement finetuning (Chapter~15); multi-agent systems
(Chapter~16); and a survey of open problems and emerging directions (Chapter~17).

Each chapter closes with a set of exercises that range from straightforward verifications to
more involved proofs and design questions.  We recommend attempting at least a few exercises
per chapter: reinforcement learning is, appropriately, a subject best understood by doing.

\medskip

\textbf{Companion code and notebooks.}\quad
This book is accompanied by a public GitHub repository containing practical Jupyter
notebooks, one for each chapter.  These notebooks provide runnable implementations of
every major algorithm discussed in the text\,---\,from $\varepsilon$-greedy bandits and
Q-learning to DQN, PPO, DPO, and GRPO\,---\,together with visualisations and guided
experiments.  All notebooks are designed to run on free-tier platforms such as Google
Colab or Kaggle, requiring no specialised hardware.  The repository also includes
detailed exercise solutions with full proofs and derivations, as well as a framework
for community translations into other languages.

The companion repository is available at:

\medskip
\centerline{\url{https://github.com/pyshka501/rl-textbook}}
\medskip

\noindent
We encourage readers to work through the notebooks alongside the relevant chapters.
The combination of mathematical rigour in the text and hands-on experimentation in
the code is, in our experience, the most effective way to develop genuine fluency in
reinforcement learning.

\medskip

We hope this book serves you well, wherever on the path from classical theory to modern
alignment practice you happen to be.

\bigskip

\hfill\textit{Moscow, 2026}
